\section{AI for Simulation (Prof.\ Pirker)}

Prof.\ Pirker is the supervisor and taught \emph{Numerical Methods in Fluid Mechanics}. Expect him
to stay close to recurrence CFD and to the figures of the thesis: he points at a figure and asks
what it shows and why. This part has three blocks: rCFD as taught in the lecture, the figures of the
thesis, and the CFD basics that rCFD stands on.

\subsection{Recurrence CFD as taught in the lecture}

\begin{qabox}{Why rCFD at all? Connect it to how Navier--Stokes is solved.}
\core{rCFD makes a bold prediction, reusing recorded flow, and then corrects it with physics, just as SIMPLE guesses a pressure and corrects it.}
\explain Many industrial flows run for a long time and show recurring patterns: bubble columns,
fluidized beds, turbulent vessels. Solving momentum and pressure for all that time is wasteful once
the flow repeats itself. The correction then fixes what the reuse breaks, for example the mass
balance.
\summary{Bold prediction + physics-based correction.}
\src{Pirker, rCFD lecture, slides 16--18.}
\end{qabox}

\begin{qabox}{Flow-based versus transport-based rCFD.}
\core{Flow-based rCFD stitches recorded flows and still solves a transport equation; transport-based rCFD replaces the equation by cell shifts and face swaps.}
\explain Flow-based rCFD (Lichtenegger and Pirker 2016) builds a long artificial flow from recorded
sequences, about 100 times faster. Transport-based rCFD (Pirker and Lichtenegger 2018) solves no
transport equation at all. My thesis uses the transport-based form, plus the correction hook and
the conservation constraint.
\nums{Steel ladle: 26~h of full CFD against 1.6~s of transport-based rCFD.}
\summary{Flow-based: stitch flows, solve transport. Transport-based: shifts and swaps.}
\src{Pirker, rCFD lecture, slides 19--34.}
\end{qabox}

\begin{qabox}{What are the recurrence norm, matrix, plateau and prominence?}
\core{The norm measures how different two states are, and the prominence measures how deep the flow dips back towards earlier states.}
\explain The norm for all pairs gives the matrix. One row of it starts at zero on the diagonal,
rises and then oscillates about a plateau. Deep dips below the plateau mean the flow comes back to
similar states, so the replay can jump there. In my tank the dips are shallow, so the replay plays
the frames in order.
\nums{The method asks for a prominence of 0.05 to 0.1. My tank: median 0.028.}
\derive
\begin{align*}
R_{fg} &= \bigl\lVert\Phi^{(f)} - \Phi^{(g)}\bigr\rVert \why{norm between two frames} \\
d(f,g) &= R_{fg}/\max R \why{normalised distance} \\
P_{\mathrm{rel}} &= \bigl\langle d_{\max} - d_{\min}\bigr\rangle / \bar{d} \why{dip depth over the plateau}
\end{align*}
\summary{Norm, matrix, row, plateau, prominence. Mine: 0.028, frames in order.}
\src{Pirker, rCFD lecture, slides 50--51; thesis, Section 2.4.}
\end{qabox}

\begin{qabox}{The three ``bold failures'' of rCFD and their corrections.}
\core{Mass imbalance, unphysical diffusion and too little recurrence, fixed by a global balance, local swaps and recurrence islands.}
\explain First, cell shifts do not conserve mass locally, so a global monitor corrects the total
with mass-acting face swaps. Second, a fixed number of swaps gives the same mixing everywhere, so
the swaps should follow the local mixing. Third, in large domains frames that match globally do not
match locally, so the domain is split into islands with their own matrices. My thesis touches all
three: the conservation constraint, the recorded diffusivity, and a flow with too little
prominence.
\summary{Imbalance $\to$ global balance. Diffusion $\to$ local swaps. Low recurrence $\to$ islands.}
\src{Pirker, rCFD lecture, slides 35--54; thesis, Sections 2.6--2.8.}
\end{qabox}

\begin{qabox}{When does rCFD work, and when does it fail?}
\core{It works when the motion repeats and the scalar does not change it, and it fails when the scalar drives the flow or the operating state changes.}
\explain In a new state the recording does not contain the motion that this state would produce.
That is my result: under the cooled lid the flow keeps renewing itself with cold plumes, a motion
that the adiabatic recording never saw. The answer of the lecture is multi-state rCFD: one database
per process state.
\summary{A database belongs to a flow regime, not to a case.}
\src{Pirker, rCFD lecture, slide 59; thesis, Section 6.4, Figure 6.6.}
\end{qabox}

\begin{qabox}{rCFD and digital twins. How would you use your work in a plant?}
\core{rCFD can run alongside the process, if live data keep aligning it.}
\explain A tank model with one mean value is instant but ignores the inside of the plant. Full CFD
shows the inside but takes days. rCFD sits in between. Run blindly forward, it collects error, so
live data should correct it. My thesis shows what such a twin must watch: a boundary condition can
move the process into a regime the database does not cover.
\nums{Lecture example: tank model under 1~s, experiment 750~s, CFD 2 days, rCFD 30~s.}
\summary{Fast but blind, exact but slow, rCFD in between. Align with live data.}
\src{Pirker, rCFD lecture, slides 59--72.}
\end{qabox}

\begin{qabox}{Is rCFD machine learning?}
\core{It is data-based simulation: an expensive ground truth makes data, a cheap model reuses it, and the question is how far it generalizes.}
\explain There is no gradient training. Choosing the next frame is a nearest-neighbour search. The
replay is a sequence of sparse linear operators, close to a learned propagator (Lumetzberger et
al.\ 2025). Possible AI steps: learn the map for unseen states, learn the missing flux as a
function of the state, or train a neural operator on the database.
\summary{Ground truth $\to$ data $\to$ cheap reuse. Generalization = which regime the database covers.}
\src{Pirker, rCFD lecture, slide 69; thesis, Section 2.2.}
\end{qabox}

\subsection{The figures of the thesis}

For every figure, have three sentences ready: what is plotted, what we see, and why.

\begin{qabox}{Figure 1.1 -- passive and non-passive transport.}
\core{A dye only rides the flow, but temperature makes the flow, and the replay cuts this feedback.}
\explain (a) A dye in the stored velocity of the adiabatic tank is only stretched. (b) In the cooled
tank, cold plumes from the lid create the motion. (c) The replay moves temperature with recorded
maps and never solves the velocity.
\summary{The whole thesis question in one figure.}
\end{qabox}

\begin{qabox}{Figure 2.2 -- recurrence of the database.}
\core{The flow has two phases, filling and settling, and no state comes back.}
\explain The distance matrix has two blocks split at frame 315. The rows rise to a plateau and none
returns towards zero. The prominence stays below what the method needs, so the replay plays the
frames in order.
\nums{Median prominence 0.028. Nearest-neighbour distance 0.013, only because neighbouring frames are similar.}
\summary{Two blocks, no return, frames in order.}
\end{qabox}

\begin{qabox}{Figure 3.3 -- vertical heat transport and plume scale.}
\core{Under the cooled lid, the flow carries heat about 600 times more than conduction.}
\explain The plane-mean vertical heat flux is split into the part carried by the flow and the part
carried by conduction. The temperature departures under the lid have the size of whole plumes. So
cooling is a transport process, and a diffusive step cannot copy it.
\summary{Cooling is transport, not diffusion.}
\end{qabox}

\begin{qabox}{Figures 3.4 and 3.5 -- the two regimes.}
\core{One boundary condition switches the flow from one regime to another.}
\explain With the adiabatic lid the jet rises, spreads under the lid, and the warm layer then rests.
With the cooled lid the layer above the front overturns in a row of cells driven by cold plumes.
Same geometry, same inflow, two different flows.
\summary{Same tank, same inflow, two flows.}
\end{qabox}

\begin{qabox}{Figure 5.1 -- the standard replay over the face-exchange sweep.}
\core{More mixing heats the tank faster during the inflow and drains it faster afterwards.}
\explain The mean temperature and the height of the heat are shown for seven strengths. The height
is best at intermediate strengths, the mean temperature with little mixing. So no single strength is
best for both.
\summary{No single best $\fsd$.}
\end{qabox}

\begin{qabox}{Figure 5.2 -- the adiabatic tank through the run.}
\core{The recorded maps keep the front thin, the face exchange smears it, and the correction makes it thin again.}
\explain The mid-plane temperature is shown at four times, with the front width in each panel. The
replay with maps only keeps the front as thin as the reference. So the smearing comes from the
exchange, not from the maps.
\summary{Maps: thin. Exchange: smeared. Correction: thin again.}
\end{qabox}

\begin{qabox}{Figure 5.3 -- space--time diagram.}
\core{It shows the warm layer growing downwards and where the standard replay is too warm or too cold.}
\explain The horizontal mean temperature is plotted against height and time. The difference panel
shows the error of the standard replay, and the thickness panel shows that its front stays too
thick. The corrected replay follows the reference.
\summary{Height against time: the layer grows, the error is visible.}
\end{qabox}

\begin{qabox}{Figure 5.6 -- every correction run in one map.}
\core{Only the vertically weighted form brings both errors down while the front reaches the reference width.}
\explain Each marker is one replay: trace error and end-profile error against the front width. The
other forms reach a thin front only with a large error. The working range of the vertical form is
$s = 0.40$ to 0.50.
\nums{Reference front 5.0~mm.}
\summary{Direction matters more than strength.}
\end{qabox}

\begin{qabox}{Figure 5.7 -- buoyancy frequency across the front.}
\core{A sharper front is a stronger barrier to vertical motion.}
\explain The buoyancy frequency measures how strongly the stratification resists vertical motion.
The standard replay holds a barrier about half as strong as the reference. The corrected replay
reaches the reference value after the inflow.
\derive
\begin{align*}
N^2 &= g\,\beta\,\frac{\partial T}{\partial y} \why{buoyancy frequency, Eq.~(3.3)}
\end{align*}
\summary{Steeper gradient, larger $N$, stronger barrier.}
\end{qabox}

\begin{qabox}{Figure 5.8 -- end profiles through the front.}
\core{Read the width always together with the profile.}
\explain Without correction the front is wide and too high. The unweighted correction makes it thin
but puts it in the wrong place. The vertical weight follows the reference. A saw-toothed front can
have the right width and the wrong shape.
\summary{Right width is not enough: check the shape.}
\end{qabox}

\begin{qabox}{Figure 5.10 -- outlet temperature from the energy balance.}
\core{The outlet sits just below the front, so a smeared front turns into a heat leak.}
\explain While the inflow is warm the standard replay keeps heat in the tank, and afterwards it
drains it. The vertical correction ends much closer. So a numerical error at the front becomes an
error in the energy budget.
\nums{Standard: $+2.7$~K kept at $t = 315$~s, ends at $-1.9$~K. Vertical correction ends at
$-0.5$~K.}
\derive
\begin{align*}
M\frac{\mathrm{d}\bar T}{\mathrm{d}t} &= \dot m\,(T_{\mathrm{in}} - T_{\mathrm{out}}) \why{energy in minus out, Eq.~(4.10)} \\
T_{\mathrm{out}} &= T_{\mathrm{in}} - \tau\,\frac{\mathrm{d}\bar T}{\mathrm{d}t},\quad \tau = M/\dot m \why{solve for the outlet, Eq.~(4.11)}
\end{align*}
\summary{Smeared front + outlet below it = heat leak.}
\end{qabox}

\begin{qabox}{Figure 5.11 -- state space.}
\core{The correction keeps the replay close to the states that the recording spans.}
\explain The two leading POD modes of the reference field are shown, and the distance of each run
from the space of the first ten modes. The standard replay leaves this space and does not return.
The corrected replay returns after the inflow.
\summary{Corrected replay stays on the recorded state space.}
\end{qabox}

\begin{qabox}{Figure 6.2 -- wall-cooled end profiles.}
\core{The standard replay holds nearly the right heat, but in the wrong place.}
\explain The bulk is too warm and the top layer too cold. The correction moves the heat back. This
is why I report the mean temperature and the profile together.
\nums{Standard replay: profile 2.79~K. Corrected at $s = 0.475$: 0.59~K.}
\summary{Same heat, wrong place.}
\end{qabox}

\begin{qabox}{Figure 6.3 -- wall-cooled traces.}
\core{The local correction follows the heat content but lags in layering while the tank fills.}
\explain The red band shows how far the local correction lags in heterogeneity during the warm
inflow. Only the recorded flux also follows the heterogeneity. A lag in layering is missing
transport, not missing sharpness.
\summary{Layering lag = missing transport.}
\end{qabox}

\begin{qabox}{Figures 6.4 and 6.5 -- end state and plumes.}
\core{The departure is smallest at $s = 0.55$, and only the corrected replay has plumes that reach the front.}
\explain Along the strength ladder the cold top and the warm bulk fade together. At $s = 0.8$ the
top turns warm. In 3D, the plumes of the standard replay end halfway down, and with the correction
they reach the front again. The plumes were in the recorded maps all along.
\nums{Smallest mean departure 0.22~K at $s = 0.55$.}
\summary{Right layering lets the recorded plumes show.}
\end{qabox}

\begin{qabox}{Figure 6.6 -- recurrence of the two recordings.}
\core{After the inflow the cooled flow keeps renewing itself, and the adiabatic flow hardly changes.}
\explain At short lags the cooled flow changes about ten times faster. So the cooled case is a
different regime, and a database recorded without cooling cannot contain its motion.
\summary{Different regime $\to$ different database.}
\end{qabox}

\subsection{CFD basics behind rCFD}

\begin{qabox}{Reynolds number and other dimensionless numbers.}
\core{The Reynolds number compares inertia with viscous forces.}
\explain Pipe flow is laminar below about 2300 and turbulent above about 4000. Other numbers: the
P\'eclet number compares convection with diffusion, the Prandtl number momentum diffusion with heat
diffusion, the Froude number inertia with gravity, and the Richardson number buoyancy with inertia.
\nums{My case: $Re \approx 99$, laminar. $Ri = 1.84$.}
\derive
\begin{align*}
Re &= \frac{UL}{\nu},\quad Pe = \frac{UL}{\Gamma},\quad Pr = \frac{\nu}{\alpha},\quad Ri = \frac{g'L}{U^2} \why{ratios of effects}
\end{align*}
\summary{Re: inertia/viscosity. Pe: convection/diffusion. Ri: buoyancy/inertia.}
\src{Pirker, introduction slides; thesis, Section 3.3.}
\end{qabox}

\begin{qabox}{Governing equations and the material derivative.}
\core{Mass, momentum and energy are conserved, and the material derivative links the moving parcel with the fixed control volume.}
\explain The Lagrangian view follows a fluid parcel. The Eulerian view looks at a fixed point. The
material derivative is the rate of change seen by the moving parcel.
\derive
\begin{align*}
\frac{\mathrm{d}}{\mathrm{d}t}\phi\bigl(\mathbf{x}(t),t\bigr) &= \frac{\partial\phi}{\partial t} + \sum_i\frac{\partial\phi}{\partial x_i}\frac{\mathrm{d}x_i}{\mathrm{d}t} \why{chain rule along the path} \\
&= \frac{\partial\phi}{\partial t} + \mathbf{u}\cdot\nabla\phi = \frac{\mathrm{D}\phi}{\mathrm{D}t} \why{the parcel moves with $\mathrm{d}\mathbf{x}/\mathrm{d}t = \mathbf{u}$} \\
\frac{\partial\rho}{\partial t} + \nabla\cdot(\rho\mathbf{u}) &= 0 \why{mass} \\
\frac{\partial(\rho\mathbf{u})}{\partial t} + \nabla\cdot(\rho\mathbf{u}\mathbf{u}) &= -\nabla p + \nabla\cdot\bm\tau + \rho\mathbf{g} \why{momentum}
\end{align*}
\summary{Change in time + convection = pressure + viscous + body force.}
\src{Pirker, governing-equation slides.}
\end{qabox}

\begin{qabox}{The finite-volume method.}
\core{Integrate the equation over each cell, and the divergence theorem turns everything into fluxes through faces.}
\explain What leaves one cell enters its neighbour, so the method is conservative. Each cell value
is linked to its neighbours by one linear equation. If all coefficients are positive, the new value
is a weighted mean, and no new extremum can appear. The same rule protects the face exchange in
rCFD.
\derive
\begin{align*}
\frac{\partial c}{\partial t} + \nabla\cdot(\mathbf{u}c) &= \nabla\cdot(\Gamma\nabla c) \why{transport equation} \\
V\frac{\mathrm{d}c_P}{\mathrm{d}t} + \oint\mathbf{u}c\cdot\mathbf{n}\,\mathrm{d}A &= \oint\Gamma\,\nabla c\cdot\mathbf{n}\,\mathrm{d}A \why{integrate, divergence theorem} \\
a_P\,c_P &= \sum_{NB} a_{NB}\,c_{NB} + b \why{fluxes per face: the computational molecule}
\end{align*}
\summary{Integrate, fluxes through faces, conservative, positive coefficients.}
\src{Pirker, finite-volume slides.}
\end{qabox}

\begin{qabox}{Discretization schemes and numerical diffusion.}
\core{First-order upwind is stable but adds false diffusion; central differencing is more accurate but can oscillate.}
\explain Upwind takes the value from upstream. Its truncation error acts like an extra diffusivity
that smears sharp gradients. Central differencing oscillates when the cell P\'eclet number is above
2. Finer or flow-aligned grids and higher-order schemes reduce numerical diffusion. My thesis is a
numerical-diffusion story: the face exchange smears the front, and the correction undoes it.
\derive
\begin{align*}
c_{i-1} &= c_i - \Delta x\,\partial_x c + \tfrac12\Delta x^2\,\partial_{xx}c - \dots \why{Taylor series} \\
\frac{c_i - c_{i-1}}{\Delta x} &= \partial_x c - \tfrac12\Delta x\,\partial_{xx}c + \dots \why{upwind difference, $u > 0$} \\
\partial_t c + u\,\partial_x c &= \tfrac12 u\Delta x\,\partial_{xx}c + \dots \why{what upwind really solves} \\
\Gamma_{\mathrm{num}} &= \tfrac12\,u\,\Delta x \why{numerical diffusivity}
\end{align*}
\summary{Upwind: stable, diffusive, $\Gamma_{\mathrm{num}} = u\Delta x/2$. Central: accurate, oscillates if $Pe > 2$.}
\src{Pirker, finite-volume slides II.}
\end{qabox}

\begin{qabox}{Explicit versus implicit time stepping; Courant and diffusion numbers.}
\core{Explicit steps are cheap but limited by stability; implicit steps solve a system but are unconditionally stable.}
\explain Explicit convection needs a Courant number below one: information must not travel more
than one cell per step. Explicit diffusion needs a diffusion number below one half. This is exactly
the cap that stopped the recorded diffusivity in my replay.
\derive
\begin{align*}
c_i^{n+1} &= c_i^n + d\,(c_{i+1}^n - 2c_i^n + c_{i-1}^n),\quad d = \frac{\Gamma\,\Delta t}{\Delta x^2} \why{explicit diffusion} \\
G &= 1 - 4d\,\sin^2(k\Delta x/2) \why{insert a Fourier mode} \\
\lvert G\rvert \le 1 &\iff 4d \le 2 \iff d \le \tfrac12 \why{stable for every mode}
\end{align*}
\nums{$Co = u\Delta t/\Delta x < 1$. $d < 0.5$.}
\summary{$Co < 1$, $d < 0.5$. My cap is the explicit diffusion limit.}
\src{Pirker, finite-volume slides; thesis, Section 5.2.1.}
\end{qabox}

\begin{qabox}{The SIMPLE algorithm.}
\core{Guess the pressure, solve the velocity, correct the pressure from continuity, and repeat.}
\explain SIMPLE stands for Semi-Implicit Method for Pressure-Linked Equations (Patankar). The
velocity correction depends on the pressure correction. Putting it into continuity gives an
equation for the pressure correction. Then both fields are corrected, with under-relaxation, until
the residuals are small.
\derive
\begin{align*}
&\mathbf{u}^{*} \leftarrow \text{momentum with } p^{*} \why{1: predict the velocity} \\
&u' = d\,(p'_P - p'_E) \why{2: velocity correction from pressure correction} \\
&\nabla\cdot(\mathbf{u}^{*} + \mathbf{u}') = 0 \Rightarrow p' \why{3: continuity gives $p'$} \\
&p = p^{*} + \alpha_p\,p',\quad \mathbf{u} = \mathbf{u}^{*} + \mathbf{u}' \why{4: correct, then repeat}
\end{align*}
\summary{Guess $p$, solve $u$, correct from continuity, repeat.}
\src{Pirker, finite-volume slides.}
\end{qabox}

\begin{qabox}{Turbulence: phenomena, Reynolds decomposition, models.}
\core{Turbulence is chaotic and mixes strongly; RANS models all eddies, LES models only the small ones, DNS models none.}
\explain Energy enters at the large eddies and cascades down to the Kolmogorov scale, where
viscosity dissipates it. RANS splits the velocity into a mean and a fluctuation, and the Reynolds
stresses need a model, usually an eddy viscosity. My case needs none of this, because it is
laminar.
\derive
\begin{align*}
u &= \bar u + u' \why{Reynolds decomposition} \\
\nu_t &= C_\mu\,k^2/\varepsilon \why{$k$--$\varepsilon$ eddy viscosity} \\
\eta &= \nu^{a}\varepsilon^{b}:\ \ \mathrm{m} = (\mathrm{m^2/s})^{a}(\mathrm{m^2/s^3})^{b} \why{Kolmogorov scale from $\nu$ and $\varepsilon$ only} \\
&\Rightarrow\ 2a + 2b = 1,\ \ a + 3b = 0 \why{match metres and seconds} \\
\eta &= (\nu^3/\varepsilon)^{1/4} \why{$a = 3/4$, $b = -1/4$}
\end{align*}
\summary{Cascade to $\eta$. RANS all, LES small, DNS none.}
\src{Pirker, turbulence slides.}
\end{qabox}

\begin{qabox}{Wall boundary layers, $y^+$ and wall functions.}
\core{Near a wall the velocity drops to zero, and $y^+$ says how far the first cell sits from the wall in viscous units.}
\explain The viscous sublayer lies below $y^+ \approx 5$ and the log layer above about 30. Wall
functions place the first cell in the log layer and model the rest, which is cheap. A wall-resolved
simulation needs $y^+ \approx 1$ at the first cell.
\derive
\begin{align*}
u_\tau &= \sqrt{\tau_w/\rho},\quad y^+ = \frac{u_\tau\,y}{\nu} \why{friction velocity, wall distance in viscous units}
\end{align*}
\summary{$y^+ < 5$ viscous, $> 30$ log layer. Wall functions: first cell in the log layer.}
\src{Pirker, turbulence slides III.}
\end{qabox}

\begin{qabox}{Species and heat transport; the Boussinesq approximation.}
\core{Temperature and species follow the same transport equation, and Boussinesq keeps the density constant except where gravity acts.}
\explain If the density depends on temperature, flow and temperature are coupled. For small density
changes we keep the density constant everywhere except in the gravity term. In my case the change
is 1.5 per cent.
\derive
\begin{align*}
\frac{\partial c}{\partial t} + \nabla\cdot(\mathbf{u}c) &= \nabla\cdot(\Gamma\nabla c) + S \why{scalar transport} \\
\rho\,\mathbf{g} &\approx \rho_0\,\mathbf{g}\bigl(1 - \beta(T-T_0)\bigr) \why{Boussinesq: density only in the buoyancy term}
\end{align*}
\summary{Constant density, except where gravity sees it.}
\src{Pirker, species and heat slides; thesis, Section 2.1.}
\end{qabox}

\begin{qabox}{Multiphase flows: Euler--Euler versus Euler--Lagrange.}
\core{Euler--Euler treats each phase as a continuum; Euler--Lagrange tracks every particle.}
\explain Euler--Euler uses volume fractions and suits dense systems at lower cost. Euler--Lagrange
follows particles in the fluid, and with particle contacts it becomes CFD--DEM: accurate but
expensive. rCFD was first used for bubble columns and fluidized beds.
\summary{Continuum phases versus tracked particles.}
\src{Pirker, multiphase slides.}
\end{qabox}
